\section{Monomial matrices and their common fixed spaces}
\label{sec:monomial_matrices}

\begin{lemma}[Matrix units span all matrices]
    \label{thm:monomial_matrix_units_span}
    \lean{Submodule.eq_top_of_forall_single_mem}
    \leanok
    Let $R$ be a semiring and let $m,n$ be finite index sets. A submodule of
    the $m\times n$ matrices over $R$ that contains every matrix unit
    $E_{ij}$ is the whole space.
\end{lemma}

\begin{proof}\leanok
    Every matrix is the sum $X=\sum_{i,j}X_{ij}E_{ij}$, and a submodule is
    closed under scalar multiples and finite sums.
\end{proof}

\begin{definition}[Monomial matrix]
    \label{def:mpug_monomial_matrix}
    \lean{Matrix.monomial,
        Matrix.monomial_apply,
        Matrix.monomial_apply_perm,
        Matrix.monomial_apply_of_ne}
    \leanok
    Let $\sigma$ be a permutation of a finite index set $S$ and let
    $\varphi\colon S\to\C$. The monomial matrix $M(\sigma,\varphi)$ is the
    operator on $\C^S$ with
    \begin{align}
        M(\sigma,\varphi)\ket{s} & =\varphi(s)\ket{\sigma(s)}.
        \label{eq:mpug_monomial_matrix}
    \end{align}
    Its entry in row $\sigma(s)$ and column $s$ is $\varphi(s)$, and every
    other entry vanishes.
\end{definition}

\begin{lemma}[Calculus of monomial matrices]
    \label{thm:mpug_monomial_matrix_laws}
    \lean{Matrix.monomial_mulVec,
        Matrix.monomial_mulVec_single,
        Matrix.monomial_mulVec_eq_self_iff,
        Matrix.monomial_mul_monomial,
        Matrix.conjTranspose_monomial,
        Matrix.smul_monomial,
        Matrix.neg_monomial,
        Matrix.monomial_one,
        Matrix.reindex_monomial,
        Matrix.monomial_mem_unitaryGroup,
        Matrix.trace_monomial,
        Matrix.trace_monomial_eq_zero}
    \leanok
    \uses{def:mpug_monomial_matrix}
    A vector $v\in\C^S$ is fixed by $M(\sigma,\varphi)$ if and only if
    $v_{\sigma(s)}=\varphi(s)v_s$ for every $s\in S$. Products, adjoints,
    scalar multiples, negatives, and simultaneous relabellings of monomial
    matrices are monomial:
    \begin{align}
        M(\sigma,\varphi)M(\tau,\psi)
         & =M\bigl(\sigma\tau,\ s\mapsto\varphi(\tau(s))\psi(s)\bigr),
        \label{eq:mpug_monomial_mul}                                               \\
        M(\sigma,\varphi)^\dagger
         & =M\bigl(\sigma^{-1},\ t\mapsto\overline{\varphi(\sigma^{-1}(t))}\bigr).
        \label{eq:mpug_monomial_adjoint}
    \end{align}
    The identity is $M(\id,1)$, and $M(\sigma,\varphi)$ is unitary when
    $|\varphi(s)|=1$ for every $s$. The trace is
    $\tr M(\sigma,\varphi)=\sum_{s:\,\sigma(s)=s}\varphi(s)$, which vanishes
    when $\sigma$ has no fixed point.
\end{lemma}

\begin{proof}\leanok
    \uses{def:mpug_monomial_matrix}
    The action of $M(\sigma,\varphi)$ on a vector is
    $(M(\sigma,\varphi)v)_t=\varphi(\sigma^{-1}(t))v_{\sigma^{-1}(t)}$, which
    gives the fixed-vector criterion after the substitution $t=\sigma(s)$.
    The remaining identities are entrywise: in a product only the column
    $\tau(s)$ of the first factor meets the row $\tau(s)$ of the second, the
    adjoint transposes the single nonzero entry of each column, and a
    relabelling conjugates the permutation. For unitarity,
    (\ref{eq:mpug_monomial_mul}) and (\ref{eq:mpug_monomial_adjoint}) give
    $M(\sigma,\varphi)^\dagger M(\sigma,\varphi)=M(\id,|\varphi|^2)=\Id$.
    The diagonal entry at $s$ is $\varphi(s)$ if $\sigma(s)=s$ and zero
    otherwise, which gives the trace.
\end{proof}

\begin{lemma}[Characters of binary configurations]
    \label{thm:monomial_binary_character_sum}
    \lean{ZMod.neg_one_pow_val_add,
        Fintype.sum_prod_neg_one_pow_val_mul}
    \leanok
    For $x,y\in\mathbb F_2$, the sign $(-1)^x$ satisfies
    $(-1)^{x+y}=(-1)^x(-1)^y$. For a finite set $I$ and $c\in\mathbb F_2^I$,
    \begin{align}
        \sum_{a\in\mathbb F_2^I}\prod_{i\in I}(-1)^{a_ic_i}
         & =\begin{cases}
                2^{|I|} & \text{if }c=0,    \\
                0       & \text{otherwise}.
            \end{cases}
        \label{eq:monomial_binary_character_sum}
    \end{align}
\end{lemma}

\begin{proof}\leanok
    The sum of a product over independent coordinates is the product of the
    one-coordinate sums $\sum_{x\in\mathbb F_2}(-1)^{xc_i}$, which equal $2$
    if $c_i=0$ and $0$ otherwise.
\end{proof}

\begin{lemma}[Flat phases on the hypercube have a potential]
    \label{thm:mpug_hypercube_potential}
    \lean{Hypercube.hypercube_induction,
        Hypercube.exists_potential_of_flat}
    \leanok
    Let $M$ be a commutative group and let
    $\varphi_j\colon\mathbb F_2^r\to M$ for $0\leq j<r$ be phases attached
    to the edges $(\gamma,\gamma+e_j)$ of the hypercube. Assume that
    \begin{align}
        \varphi_j(\gamma)\varphi_j(\gamma+e_j)
         & =1,
        \label{eq:mpug_hypercube_involutive}       \\
        \varphi_j(\gamma)\varphi_k(\gamma+e_j)
         & =\varphi_k(\gamma)\varphi_j(\gamma+e_k)
        \label{eq:mpug_hypercube_flat}
    \end{align}
    for all $j,k$ and $\gamma$. Then there is $\Phi\colon\mathbb F_2^r\to M$
    with $\Phi(0)=1$ and $\Phi(\gamma+e_j)=\varphi_j(\gamma)\Phi(\gamma)$ for
    every $j$ and $\gamma$. In addition, a property of vectors of
    $\mathbb F_2^r$ that holds at $0$ and is preserved by every coordinate
    flip holds everywhere.
\end{lemma}

\begin{proof}\leanok
    For the connectivity claim, every $\gamma$ is the sum of the $e_j$ over
    its support, so induction on the support carries the property from $0$ to
    $\gamma$.

    The potential is constructed by induction on $r$. Writing
    $\gamma=(a,\gamma')$ with $a\in\mathbb F_2$ and
    $\gamma'\in\mathbb F_2^{r-1}$, the phases
    $\gamma'\mapsto\varphi_{j+1}(0,\gamma')$ for $0\leq j<r-1$ satisfy
    (\ref{eq:mpug_hypercube_involutive}) and (\ref{eq:mpug_hypercube_flat})
    on $\mathbb F_2^{r-1}$, so the induction hypothesis gives a potential
    $\Phi'$ for them. Setting
    \begin{align}
        \Phi(0,\gamma') & =\Phi'(\gamma'),
        \qquad
        \Phi(1,\gamma')=\varphi_0(0,\gamma')\Phi'(\gamma'),
        \label{eq:mpug_hypercube_potential_step}
    \end{align}
    the edge relation in direction $0$ at $(0,\gamma')$ holds by definition,
    and at $(1,\gamma')$ it follows from
    (\ref{eq:mpug_hypercube_involutive}). The edge relation in a direction
    $j+1$ at $(0,\gamma')$ is the relation for $\Phi'$, and at $(1,\gamma')$
    it is the relation for $\Phi'$ combined with
    (\ref{eq:mpug_hypercube_flat}) for the square spanned by the directions
    $0$ and $j+1$ at $(0,\gamma')$.
\end{proof}

\begin{definition}[Common fixed subspace]
    \label{def:mpug_common_fixed_subspace}
    \lean{LinearMap.commonFixedSubmodule}
    \leanok
    Let $\mathcal H$ be a complex vector space and let $(P_j)_{j\in J}$ be any
    family of endomorphisms of $\mathcal H$. Their common fixed subspace is
    \begin{align}
        \mathcal V(P)
         & =\bigcap_{j\in J}\ker(P_j-\Id_{\mathcal H}).
    \end{align}
    No commutativity of the operators is required.
\end{definition}

\begin{lemma}[Membership in the common fixed subspace]
    \label{thm:mpug_mem_common_fixed_subspace}
    \lean{LinearMap.mem_commonFixedSubmodule_iff}
    \leanok
    \uses{def:mpug_common_fixed_subspace}
    A vector $\psi\in\mathcal H$ belongs to $\mathcal V(P)$ if and only if
    $P_j\psi=\psi$ for every $j\in J$.
\end{lemma}

\begin{proof}\leanok
    \uses{def:mpug_common_fixed_subspace}
    Membership in an indexed intersection is equivalent to membership in
    every factor, and membership in $\ker(P_j-\Id_{\mathcal H})$ is equivalent
    to $P_j\psi=\psi$.
\end{proof}

\begin{lemma}[Nonzero-vector lower bound]
    \label{thm:mpug_common_fixed_subspace_finrank}
    \lean{LinearMap.one_le_finrank_commonFixedSubmodule}
    \leanok
    \uses{def:mpug_common_fixed_subspace}
    If $\mathcal H$ is finite-dimensional, $\psi\in\mathcal V(P)$, and
    $\psi\neq0$, then $1\leq\dim_{\C}\mathcal V(P)$.
\end{lemma}

\begin{proof}\leanok
    The common fixed subspace is nonzero because it contains $\psi$, so its
    dimension is at least one.
\end{proof}

\begin{definition}[Trivial holonomy on a fiber]
    \label{def:mpug_trivial_holonomy}
    \lean{Hypercube.fiberFlip,
        Hypercube.fiberFlip_apply,
        Hypercube.IsTrivialHolonomy}
    \leanok
    Let $X$ be a finite set and consider the basis $X\times\mathbb F_2^r$.
    For $0\leq j<r$ let $\sigma_j$ be the coordinate flip
    $(x,\gamma)\mapsto(x,\gamma+e_j)$ and let
    $\varphi_j\colon X\times\mathbb F_2^r\to\C$. The phases have trivial
    holonomy on the fiber over $x\in X$ when
    \begin{align}
        \varphi_j(x,\gamma)\varphi_k(x,\gamma+e_j)
         & =\varphi_k(x,\gamma)\varphi_j(x,\gamma+e_k)
        \label{eq:mpug_trivial_holonomy}
    \end{align}
    for all $j,k$ and $\gamma\in\mathbb F_2^r$.
\end{definition}

\begin{lemma}[Monomial fixed-space criterion]
    \label{thm:mpug_monomial_fixed_space_criterion}
    \lean{Matrix.commonFixedSubmodule_reindex,
        Matrix.finrank_commonFixedSubmodule_reindex,
        Hypercube.add_self_eq_zero_pi_zmod_two,
        Hypercube.eq_zero_of_mem_commonFixedSubmodule_monomial_fiberFlip,
        Hypercube.finrank_commonFixedSubmodule_monomial_fiberFlip}
    \leanok
    \uses{def:mpug_monomial_matrix, def:mpug_trivial_holonomy,
        def:mpug_common_fixed_subspace}
    In the setting of Definition~\ref{def:mpug_trivial_holonomy}, let
    $T_j=M(\sigma_j,\varphi_j)$ and assume that every $T_j$ is an involution,
    that is, $\varphi_j(s)\varphi_j(\sigma_j(s))=1$ for every basis element
    $s$. Then
    \begin{align}
        \dim\mathcal V(T)
         & =\#\{x\in X:\text{the phases have trivial holonomy on the fiber over }x\}.
        \label{eq:mpug_monomial_fixed_space_count}
    \end{align}
    The dimension of a common fixed subspace is also unchanged by a
    simultaneous relabelling of the basis.
\end{lemma}

\begin{proof}\leanok
    \uses{thm:mpug_monomial_matrix_laws, thm:mpug_hypercube_potential,
        thm:mpug_mem_common_fixed_subspace}
    By Lemma~\ref{thm:mpug_monomial_matrix_laws}, a vector $v$ lies in
    $\mathcal V(T)$ exactly when
    $v_{(x,\gamma+e_j)}=\varphi_j(x,\gamma)v_{(x,\gamma)}$ for all $j$, $x$,
    and $\gamma$. Let $X_0$ be the set of fibers with trivial holonomy and
    consider the linear map $E\colon\mathcal V(T)\to\C^{X_0}$,
    $v\mapsto(v_{(x,0)})_{x\in X_0}$.

    The map $E$ is injective: if $v\in\mathcal V(T)$ vanishes at one point
    of a fiber, the edge relation and the connectivity statement of
    Lemma~\ref{thm:mpug_hypercube_potential} make it vanish on the whole
    fiber. On a fiber with trivial holonomy, $v$ therefore vanishes as soon
    as $v_{(x,0)}=0$. On a fiber where (\ref{eq:mpug_trivial_holonomy})
    fails at some $(j,k,\gamma_0)$, the two paths from $(x,\gamma_0)$ to
    $(x,\gamma_0+e_j+e_k)$ give
    \begin{align}
        v_{(x,\gamma_0+e_j+e_k)} & =\varphi_j(x,\gamma_0)\varphi_k(x,\gamma_0+e_j)\,v_{(x,\gamma_0)}, \\
        v_{(x,\gamma_0+e_j+e_k)} & =\varphi_k(x,\gamma_0)\varphi_j(x,\gamma_0+e_k)\,v_{(x,\gamma_0)}.
    \end{align}
    The two phase products differ, so $v_{(x,\gamma_0)}=0$ and again $v$
    vanishes on the fiber.

    The map $E$ is also surjective: given $w\in\C^{X_0}$,
    Lemma~\ref{thm:mpug_hypercube_potential} provides a potential $\Phi_x$
    on each fiber with trivial holonomy for the invertible phases
    $\gamma\mapsto\varphi_j(x,\gamma)$, and setting
    $v_{(x,\gamma)}=w_x\Phi_x(\gamma)$ there and $v=0$ on the other fibers
    satisfies the potential edge relation, so $v\in\mathcal V(T)$,
    and $\Phi_x(0)=1$ gives $Ev=w$.

    Hence $\dim\mathcal V(T)=\dim\C^{X_0}=\#X_0$. For the relabelling
    statement, a relabelled matrix acts by relabelling the vector, acting,
    and relabelling back, so the common fixed subspace is the image of the
    original one under the coordinate relabelling.
\end{proof}
